\section{The Second Paris Regency, 1268/69--1272}

The order sent Thomas back to Paris for an almost unprecedented second
regency, and into a triple war. On one flank the secular masters had
renewed the assault on the mendicants' place in the university, and
Thomas wrote his later defenses of the religious state. On the other
flank stood the radical Aristotelians of the arts faculty---the
``Averroists'' of the later textbooks, Siger of Brabant at their
head---whose reading of Aristotle through his great Arabic commentator
yielded conclusions no Christian could hold: one shared intellect for all
men, a world necessarily eternal. Thomas fought them with their own
weapons. His treatise \work{De unitate intellectus contra Averroistas}
(c.~1270) argues on strictly philosophical grounds that the intellect is
each man's own---for if thinking were done for us by another substance,
``this man'' would not think, merit, or be responsible at all---and ends
with a rare flash of personal heat, challenging his opponent to answer
in writing rather than to boys in corners. In December 1270 the bishop
of Paris, Stephen (\'Etienne) Tempier, condemned a first list of
propositions drawn from the radical arts teaching (the Stanford
Encyclopedia's account of the Paris condemnations records censures by
Tempier in 1270 and 1277); Thomas's own position was harder to place,
for he was fighting on a third front at once: against conservative
theologians for whom the Aristotelian instrument itself---one substantial
form in man, a philosophically defensible account of nature's
integrity---smelled of the same fire. In the middle of the battlefield
he kept producing at a rate that still defines the standard chronologies'
problems: the Second Part of the \work{Summa}, commentaries on Aristotle
work by work, disputed questions \work{De malo} and \work{De virtutibus},
the great Pauline and Johannine lectures. In 1272 the order called him
home to Naples to found a provincial \textit{studium} of theology, and
the Paris chapter of his life closed.
